% 제출본 전용 — 1장 끝에 붙는다 (1.4, 1.5).
% 조원 편성 · 추진 경과 · 양식 예시 목차와의 대응표.

\subsection{조원 편성}

\begin{table}[htbp]
  \centering
  \caption{조원 편성}
  \label{tab:v2-team}
  \begin{tabular}{llp{8.6cm}}
    \toprule
    \textbf{구분} & \textbf{이름} & \textbf{역할} \\
    \midrule
    조장 (1인) & 이찬희 &
      자료 수집·정제, 공고 원문 판독, 식별 전략 설계, 추정과 강건성 검정,
      파이프라인·검증 장치 구현, 문서 작성 전부 \\
    \bottomrule
  \end{tabular}
\end{table}

\textbf{1인 작품이다.} 병렬 작업이 불가능하므로 확장 과제(단지 좌표를 쓴 공간
회귀불연속, 거래 주체 분석, 동별 신규 입주물량 통제)는 처음부터 범위 밖에 두고,
답할 수 있는 다섯 질문에 집중했다(1.1절).

\subsection{추진 경과}

주 단위로 무엇을 알게 되었고 설계가 어떻게 바뀌었는지를 보고서로 남겼다. 그
\pipeDocs 편 가운데 주차 보고서가 추진 경과를 그대로 보여 준다\citep{site}.

\begin{table}[htbp]
  \centering
  \caption{주차별 추진 경과}
  \label{tab:v2-schedule}
  \footnotesize
  \begin{tabular}{llp{9.6cm}}
    \toprule
    \textbf{주차} & \textbf{날짜} & \textbf{한 일과 그 결과} \\
    \midrule
    1 & 09-04 & 주제 확정, 수집·정제 파이프라인, 선행연구 재현.
      \textbf{주 결과변수를 가격에서 거래량으로 바꿨다}(0001) \\
    2 & 09-11 & 공고 원문 판독 완료(78건 중 67건 기계 판독).
      \textbf{사건 시점을 효력일로 정렬하자 $-59\%$가 $-76\%$가 됐다}(0013) \\
    3 & 09-18 & 적은 클러스터 추론 도입. 이질성의 원인을 대출에서
      \textbf{허가 면적 기준}으로 바꿨다(0019 · 0022) \\
    4 & 09-23 & 2025년 처치 역전, 파급 분해, 가격지수 둘.
      \textbf{3주차에 적어 둔 예측이 자료에서 확인됐다}(0023) \\
    5 & 10-01 & 사전추세 민감도, 음의 가중치 진단, 사양 곡선.
      \textbf{M3 분석 종료} --- 목표(11/05)보다 5주 빠르다 \\
    6 & 10-08 & 2021년 B그룹으로 외적 타당성 확인.
      \textbf{두 사건의 차이가 0과 구분되지 않는다}(0031) \\
    7 & 10-15 & 클린 환경 재현. \textbf{의존 목록의 구멍 둘을 찾았다.}
      결과 대시보드와 보고서 표 생성 \\
    8 & 10-24 & 최종 보고서 초고. 전체 점검에서
      \textbf{보고서 자신이 수치 검사 밖이었던 것}을 찾아 넣었다 \\
    9--12 & 10-30 $\sim$ 11-20 & 최종 발표, 제출본 작성, 최종 재현 \\
    \bottomrule
  \end{tabular}
\end{table}

\begin{table}[htbp]
  \centering
  \caption{마일스톤}
  \label{tab:v2-milestone}
  \begin{tabular}{llll}
    \toprule
    & \textbf{내용} & \textbf{목표} & \textbf{충족} \\
    \midrule
    M1 & 주 분석 확정 & 4주차 (09/24) & 09/18 \\
    M2 & 하이라이트 결과 & 7주차 (10/15) & 09/23 \\
    M3 & 분석 종료 & 10주차 (11/05) & \textbf{10/01} \\
    M4 & 최종본 · 재현 확인 & 12주차 (11/19) & 재현 10/15 · 초고 10/24 \\
    \bottomrule
  \end{tabular}
\end{table}

\textbf{M3 를 5주 앞당긴 여유를 외적 타당성과 재현 장치에 썼다.} 6주차의 2021년
B그룹 분석(9.6절)과 7주차의 클린 환경 재현(3.6절)은 애초 계획에 없던 작업이고,
일정이 빠듯했다면 하지 못했을 것들이다.

\subsection{이 보고서와 양식 예시 목차의 대응}

목차는 지도교수 확인을 받아 자유롭게 정했다. 연구의 서사를 살리기 위해 양식
예시의 8장 구조 대신 다음과 같이 배치했으며, 양식이 요구하는 항목은 모두
들어 있다.

\begin{table}[htbp]
  \centering
  \caption{양식 예시 목차와 이 보고서의 대응}
  \label{tab:v2-map}
  \footnotesize
  \begin{tabular}{ll}
    \toprule
    \textbf{양식 예시} & \textbf{이 보고서} \\
    \midrule
    1.1 배경 / 1.2 목표 & 1.1 동기와 연구 질문 · 1.2 기여 사항 \\
    1.3 관련 조사(연구) & \textbf{2장} 제도와 선행연구 (세 편 직접 재현) \\
    2.1 전체 시스템 구조 & \textbf{3장} 자료와 시스템 (그림 \ref{fig:pipeline}) \\
    2.2 설계 및 제작 & \textbf{3.8절} 설계 구성요소 연관성 \\
    2.3 시험 및 평가 & \textbf{9장} 강건성 · \textbf{9.8절} 연관성 표 \\
    2.4 현실적 제약 조건 & \textbf{10.5절} \\
    2.5 조원 편성 / 2.6 추진 계획 & 1.3절 · 1.4절 \\
    3. 개발 및 운영 환경 & \textbf{부록 C} \\
    4. 구현 및 피드백 & \textbf{5$\sim$8장} 결과 · \textbf{9장} 보완 \\
    5. 결과와 효과 & \textbf{10장} 논의 · \textbf{11장} 결론 \\
    6. 개선 희망사항 & 10.4절 한계 · 11장 \\
    7. 참고 자료, URL, 문헌 & 참고 자료 절 \\
    8. 부록(소스코드) & \textbf{부록 D} \\
    \bottomrule
  \end{tabular}
\end{table}
